\section{Implementation}
\label{sec:implementation}

% Describe the data, the preprocessing, the structure of the code, and how it
% was verified (unit tests, comparison against scikit-learn/known solutions,
% etc.). Link to the repository.

% We should probably reference the first project \cite{riseth_2026} if we reuse code?

All methods were implemented from scratch in Python using NumPy~\cite{harris_2020}.
The source code, unit tests, and experiment scripts are available in the project
repository~\cite{riseth_2026a}. This section describes the datasets and preprocessing,
the organization of the implementation, and the procedures used to verify correctness.

\subsection{Data and preprocessing}\label{subsec:data}

For the regression task, we use Runge's function~\cite{runge_1901}, a classical
example for studying the bias--variance trade-off. It is defined as
\begin{equation}\label{eq:runge-function}
    f(x) = \frac{1}{1 + 25x^2}, \qquad x \in [-1,1].
\end{equation}
As observed in Project~1~\cite{riseth_2026}, polynomial regression can exhibit
substantial oscillations near the endpoints of this interval. We generate
$n=100$ input samples uniformly from $[-1,1]$, evaluate the function at these
inputs, and add i.i.d. Gaussian noise with standard deviation
$\sigma=0.1$ to obtain the observed targets.

For the classification task, we use the MNIST dataset~\cite{lecun_2010},
which contains 70,000 grayscale images of handwritten digits, each of size
$28 \times 28$ pixels, with labels corresponding to the ten classes
$0,\ldots,9$. [Describe the train/test split and any preprocessing applied
to the images, such as flattening, normalization, or label encoding.]

\subsection{Code structure}\label{subsec:code-structure}

We reuse the structure and the gradient machinery of the previous project~\cite{riseth_2026}.
The code is structured in a package \texttt{fys\_stk4155\_p2}, containing the
directories \texttt{nn}, \texttt{optimization}; which respectively hold the code for the
models and gradient descent.

Thin constructor of the FFNN for logistic regression, which returns a single layer with
softmax activation and cross entropy loss. This allows us to use this as a baseline model
when evaluating the later neural networks.

\subsection{Verification}\label{subsec:verification}
